\documentclass[10pt,a4paper]{amsart}
\usepackage[margin=19mm]{geometry}
\usepackage{amsmath,amssymb,mathtools}
\usepackage[scr=rsfs,cal=euler]{mathalfa}
\usepackage{xfrac}
\usepackage[unicode,hidelinks,bookmarks=false]{hyperref}
\newcommand{\ie}{i.e.\ }
\newcommand{\cf}{cf.\ }
\newcommand{\resp}{respectively\ }
\newcommand{\Subsection}{\subsection}
\providecommand{\nobreakdash}{\nobreak\hbox{-}}
\setlength{\emergencystretch}{2em}
\multlinegap0pt
\usepackage{amssymb}
\usepackage{mathtools}
\newtheorem{theorem}{Theorem}
\newtheorem{proposition}{Proposition}
\newtheorem{lemma}{Lemma}
\usepackage{cleveref}
\crefname{theorem}{Theorem}{Theorems}
\crefname{proposition}{Proposition}{Propositions}
\crefname{lemma}{Lemma}{Lemmas}
\newcommand{\C}{\mathbb{C}}
\newcommand{\abs}[1]{\left\lvert #1\right\rvert}
\newcommand{\cE}{\mathcal{E}}
\newcommand{\norm}[1]{\left\lVert #1\right\rVert}
\title{Convergence rates for pivoted QR and LU}
\author{Marc Aur\`ele Gilles}
\date{Source: arXiv:2607.26863v1 (CC BY 4.0)}
\begin{document}
\maketitle

\noindent\emph{Source and licence.} Convergence rates for pivoted QR and LU. Version arXiv:2607.26863v1 (CC BY 4.0), distributed by arXiv under the Creative Commons Attribution 4.0 International licence, http://creativecommons.org/licenses/by/4.0/, as recorded in the arXiv licence metadata for this exact version and shown on the abstract page https://arxiv.org/abs/2607.26863v1. Full licence notice: \texttt{licenses/arxiv-2607.26863-CC-BY-4.0.txt}.\par\smallskip

\noindent\textit{Editorial scope.} Counted unit: the article's theorem thm:analytic-determinant (source0.tex lines 1078--1109). The article's complete original proof of it is delivered with it (source0.tex lines 2205--2440, one \texttt{proof} environment of 236 lines, which the article places away from the statement). No mathematics is written, repaired or re-derived: the statement and the proof are the article's own lines, in the article's own notation, and no retained line is substituted or re-flowed.  Retained uncounted as the source-local context the proof needs: lines 1812--1868; lines 2052--2089.  Nothing was omitted from any retained span, so the package carries no omission record.  No retained line contains a citation, so the wrapper carries no bibliography and no bibliography entry.  No retained line contains a figure environment, an image-inclusion command or drawing code, so no image is delivered and the package declares no asset dependency.  Editorial formatting only, in addition: an AMS \texttt{amsart} preamble that restates the article's own declarations and macros used by the retained lines, namely the environments \texttt{theorem}, \texttt{proposition}, \texttt{lemma} on separate counters, together with the standard packages those lines need.  The retained environments print in this document as Theorem 1, Proposition 1, Lemma 1 in order of appearance, whereas the article prints the same items with the numbering of its own full document.  The complete native source of the article as distributed in the arXiv e-print for this exact version is bundled unchanged as \texttt{sources/206297/source0.tex}.
\begin{proposition}[Bounds on Chebyshev coefficients]
\label{prop:weighted-chebyshev}
The following statements hold.
\begin{enumerate}
\item Let $g:[-1,1]\to\C$ extend to a holomorphic function
$\widetilde g:\cE_\rho\to\C$, where $\rho>1$, and suppose
\begin{equation}
    \sup_{z\in\cE_\rho}\abs{\widetilde g(z)}\le M.
    \label{eq:weighted-chebyshev-extension-bound}
\end{equation}
Then there is $b=(b_m)_{m\ge0}\in\ell^2(\mathbb N)$ such that
\begin{equation}
    \sum_{m=0}^\infty\abs{b_m}^2\le M^2
    \label{eq:weighted-chebyshev-norm}
\end{equation}
and
\begin{equation}
    g(x)
    =
    \sum_{m=0}^\infty
    b_m\eta_m\rho^{-m}T_m(x),
    \qquad x\in [-1, 1].
    \label{eq:weighted-chebyshev-expansion}
\end{equation}
\item Let $g:[-1,1]^2\to\C$ extend to a holomorphic function
\[
    \widetilde g:
    \cE_{\rho_x}\times\cE_{\rho_y}\to\C,
\]
where $\rho_x,\rho_y>1$, and suppose
\begin{equation}
    \sup_{(z,w)\in\cE_{\rho_x}\times\cE_{\rho_y}}
    \abs{\widetilde g(z,w)}
    \le M.
    \label{eq:weighted-bivariate-chebyshev-extension-bound}
\end{equation}
Then there exist coefficients $c_{pq}\in\C$, $p,q\ge0$, such that
\begin{equation}
    \sum_{p=0}^{\infty} \sum_{q=0}^{\infty} \abs{c_{pq}}^2
    \le
    M^2
    \label{eq:weighted-bivariate-chebyshev-norm}
\end{equation}
and
\begin{equation}
    g(x,y)
    =
    \sum_{p,q\ge0}
    c_{pq}\eta_p\eta_q
    \rho_x^{-p}\rho_y^{-q}T_p(x)T_q(y),
    \qquad (x,y)\in [-1, 1]^2.
    \label{eq:weighted-bivariate-chebyshev-expansion}
\end{equation}
\end{enumerate}
Both series converge absolutely and uniformly on their respective
domains.
\end{proposition}

\begin{lemma}[Matrix determinant bounds]
\label{lem:analytic-finite-volume}
\leavevmode
\begin{enumerate}

\item Let \(u_0,\ldots,u_N\in\C^n\), where \(N\ge n-1\), and set
\[
    U=[u_0\ \cdots\ u_N].
\]
Then
\begin{equation}
    \det(UU^*)
    \le
    \prod_{r=0}^{n-1}
    \left(
      \sum_{m=r}^N\norm{u_m}_2^2
    \right).
    \label{eq:analytic-gram-tail}
\end{equation}

\item Let
\[
    U\in\C^{n\times p},
    \qquad
    C\in\C^{p\times q},
    \qquad
    V\in\C^{n\times q}.
\]
Then
\begin{equation}
    \abs{\det(UCV^*)}
    \le
    \left(\frac{\norm{C}_F}{\sqrt n}\right)^n
    \sqrt{\det(UU^*)\det(VV^*)}.
    \label{eq:analytic-finite-factorization}
\end{equation}
\end{enumerate}
\end{lemma}

\begin{theorem}[Analytic determinant bounds]
\label{thm:analytic-determinant}
Let $n\ge1$.
\begin{enumerate}

\item
Let $\rho>1$ and suppose that
$\mathcal A_\rho^{(x)}(f)\le M$. Then, for arbitrary
$x_1,\ldots,x_n,y_1,\ldots,y_n\in [-1, 1]$,
\begin{equation}
    \abs{\det[f(x_i,y_j)]_{i,j=1}^n}^{1/n}
    \le
    \frac{2M\sqrt n}{\sqrt{1-\rho^{-2}}}
    \rho^{-(n-1)/2}.
    \label{eq:analytic-determinant}
\end{equation}

\item
Let $\rho_x,\rho_y>1$, and suppose that
$\mathcal A_{\rho_x,\rho_y}^{(x,y)}(f)\le M$. Then, for arbitrary
$x_1,\ldots,x_n,y_1,\ldots,y_n\in [-1, 1]$,
\begin{equation}
    \abs{\det[f(x_i,y_j)]_{i,j=1}^n}^{1/n}
    \le
    \frac{4M\sqrt n}
         {\sqrt{(1-\rho_x^{-2})(1-\rho_y^{-2})}}
    (\rho_x\rho_y)^{-(n-1)/2}.
    \label{eq:mixed-analytic-determinant}
\end{equation}

\end{enumerate}
\end{theorem}

\begin{proof}[Proof of Theorem~\ref{thm:analytic-determinant}]
For the one-sided estimate, put $
    q:=\rho^{-1}.
$ For each sampled slice $
    g_j(x):=f(x,y_j),
$ 
Proposition~\ref{prop:weighted-chebyshev} gives a sequence
\(b_j\in\ell^2(\mathbb N)\) such that
$
    \norm{b_j}_{\ell^2}\le M
$
and
\[
    f(x,y_j)
    =
    \sum_{m=0}^\infty
    (b_j)_m\eta_mq^mT_m(x).
\]
For \(N\ge n-1\), define
\[
    u_m
    :=
    \eta_mq^m
    (T_m(x_1),\ldots,T_m(x_n))^T,
    \qquad
    m\ge0,
\]
and set
$
    U_N:=[u_0\ \cdots\ u_N].
$
Let \(B_N\in\C^{(N+1)\times n}\) be given by
$
    (B_N)_{m+1,j}:=(b_j)_m,
$
and set
$
    A_N:=U_NB_N.
$
By Proposition~\ref{prop:weighted-chebyshev}, for each \(i,j\),
\[
   (A_N)_{ij}
    =
    \sum_{m=0}^N
    (b_j)_m\eta_mq^mT_m(x_i)
    \longrightarrow
    f(x_i,y_j) \quad
     \text{ as } N\to\infty.
\]
Thus, with $A:=[f(x_i,y_j)]_{i,j=1}^n$, continuity of the determinant
gives $\det A_N\to\det A$.

Since
\[
    \norm{B_N}_F^2
    =
    \sum_{j=1}^n\sum_{m=0}^N\abs{(b_j)_m}^2
    \le
    nM^2,
\]
part 2 of Lemma~\ref{lem:analytic-finite-volume}, applied with
\(V=I_n\), gives
\begin{equation}
    \abs{\det A_N}
    \le
    M^n\sqrt{\det(U_NU_N^*)}.
    \label{eq:analytic-one-sided-finite}
\end{equation}

Since \(\abs{T_m(x)}\le1\) on \([-1, 1]\),
\[
    \norm{u_0}_2^2=n,
    \qquad
    \norm{u_m}_2^2\le4nq^{2m}
    \quad(m\ge1).
\]
Thus, for every \(r\ge0\),
\[
    \sum_{m=r}^\infty\norm{u_m}_2^2
    \le
    \frac{4n}{1-q^2}q^{2r}.
\]
Part 1 of Lemma~\ref{lem:analytic-finite-volume} therefore gives
\begin{equation}
\begin{aligned}
    \sqrt{\det(U_NU_N^*)}
    &\le
    \prod_{r=0}^{n-1}
    \left(
      \sum_{m=r}^\infty\norm{u_m}_2^2
    \right)^{1/2}\\
    &\le
    \left(
      \frac{2\sqrt n}{\sqrt{1-q^2}}
    \right)^n
    q^{n(n-1)/2}.
\end{aligned}
\label{eq:analytic-one-sided-volume}
\end{equation}
Combining this with \eqref{eq:analytic-one-sided-finite} and letting
\(N\to\infty\) yields
\[
    \abs{\det A}
    \le
    \left(
      \frac{2M\sqrt n}{\sqrt{1-q^2}}
    \right)^n
    q^{n(n-1)/2}.
\]
Taking \(n\)th roots and substituting \(q=\rho^{-1}\) proves
\eqref{eq:analytic-determinant}.

For the mixed estimate, put $
    q_x:=\rho_x^{-1},
    q_y:=\rho_y^{-1}.
$ 
Proposition~\ref{prop:weighted-chebyshev} gives coefficients
\((c_{pq})_{p,q\ge0}\) such that
$
    \sum_{p,q\ge0}\abs{c_{pq}}^2\le M^2
$
and
\[
    f(x,y)
    =
    \sum_{p,q\ge0}
    c_{pq}\eta_p\eta_q
    q_x^pq_y^qT_p(x)T_q(y).
\]

For \(N\ge n-1\), define
\[
    u_p
    :=
    \eta_pq_x^p
    (T_p(x_1),\ldots,T_p(x_n))^T,
    \qquad
    p\ge0,
\]
and
\[
    v_q
    :=
    \eta_qq_y^q
    (T_q(y_1),\ldots,T_q(y_n))^T,
    \qquad
    q\ge0.
\]
Set
\[
    U_N:=[u_0\ \cdots\ u_N],
    \qquad
    V_N:=[v_0\ \cdots\ v_N],
    \qquad
    C_N:=[c_{pq}]_{p,q=0}^N, 
    \qquad
    A_N:=U_NC_NV_N^*.
\]
By the absolute and uniform convergence in
Proposition~\ref{prop:weighted-chebyshev}, for each \(i,j\),
\[
    (A_N)_{ij}
    =
    \sum_{p,q=0}^N
    c_{pq}\eta_p\eta_q
    q_x^pq_y^qT_p(x_i)T_q(y_j)
    \longrightarrow
    f(x_i,y_j) \qquad
     \text{ as } N\to\infty.
\]
Hence
$
    A_N\longrightarrow
    A:=[f(x_i,y_j)]_{i,j=1}^n
$
entrywise and
$
    \det A_N\longrightarrow\det A.
$

Moreover,
\[
    \norm{C_N}_F^2
    =
    \sum_{p,q=0}^N\abs{c_{pq}}^2
    \le M^2.
\]
Part 2 of Lemma~\ref{lem:analytic-finite-volume} gives
\begin{equation}
\begin{aligned}
    \abs{\det A_N}
    \le
    \left(\frac{M}{\sqrt n}\right)^n
    \sqrt{
      \det(U_NU_N^*)\det(V_NV_N^*)
    }.
    \label{eq:analytic-mixed-finite}
\end{aligned}
\end{equation}

The same calculation as in \cref{eq:analytic-one-sided-volume} gives
\[
    \sqrt{\det(U_NU_N^*)}
    \le
    \left(
      \frac{2\sqrt n}{\sqrt{1-q_x^2}}
    \right)^n
    q_x^{n(n-1)/2},
\]
and
\[
    \sqrt{\det(V_NV_N^*)}
    \le
    \left(
      \frac{2\sqrt n}{\sqrt{1-q_y^2}}
    \right)^n
    q_y^{n(n-1)/2}.
\]
Substituting these bounds into
\eqref{eq:analytic-mixed-finite} yields
\[
    \abs{\det A_N}
    \le
    \left(
      \frac{
        4M\sqrt n
      }{
        \sqrt{(1-q_x^2)(1-q_y^2)}
      }
    \right)^n
    (q_xq_y)^{n(n-1)/2}.
\]
Letting \(N\to\infty\), taking \(n\)th roots, and substituting
\(q_x=\rho_x^{-1}\), \(q_y=\rho_y^{-1}\), proves
\eqref{eq:mixed-analytic-determinant}.
\end{proof}

\end{document}
